\chapter*{Abstract}
\addcontentsline{toc}{chapter}{Abstract}

Quantization makes a neural network cheap enough to deploy, and it is usually
validated by checking that accuracy does not fall. On convolutional networks,
accuracy has been shown to survive quantization while the explanations of the
model's decisions do not. This thesis asks the spiking analogue: \emph{does a
quantized spiking neural network base its decisions on the same evidence as its
full-precision counterpart, and at what energy cost?} A spiking network carries
a membrane potential that can be quantized independently of the weights, and its
explanations are sequences of maps; the working hypothesis is that quantizing
the membrane degrades the temporal structure of an explanation before
quantizing the weights degrades its spatial structure.

A four-block convolutional spiking network is trained on four classes of
transient noise from the LIGO Gravity Spy release, reaching $0.9833 \pm 0.0014$
validation accuracy over three seeds under hardware-realizable constraints, and
its decisions are explained with the gradient-free Spike Activation Map. Before
any comparison, the explanation metrics are calibrated: two full-precision
models differing only in seed produce maps correlating at $0.153$ on the worst
class, against a numerical floor of $0.975$. Spatial metrics are therefore valid
only between a model and its own parent, and every quantized model is
fine-tuned from, and compared with, the full-precision model of its own seed.

Weights only, membrane only, and both are quantized at 8, 4 and 2 bits. At 8
and 4 bits accuracy and explanations stay at the level of the parent in every
arm. At 2 bits the arms separate: ternary weights collapse the classifier, while
a 2-bit membrane keeps accuracy at parity ($0.982$) and yet changes the
explanation, whose highlighted region overlaps the parent's by $0.44$ against a
floor of $0.70$ --- a change the spatial correlation ($0.92$) would have reported
as a preserved explanation. The new maps remain faithful by the deletion metric:
they point at evidence the model uses, but not at the evidence the original
pointed at. The hypothesis is not supported as an ordering of explanation
metrics, but its mechanism is: the 2-bit membrane makes the spike trains of the
first layer perfectly periodic, a temporal change the explanation's temporal
metric cannot see. That configuration is also dominated on energy and memory, so
there is no trade-off between efficiency and explainability here: a joint 4-bit
model keeps both at about one sixtieth of the estimated energy and one eighth of
the memory of its parent.
